\documentclass[]{../../../../tex/classes/styledArticle}
\usepackage{../../../../tex/packages/hebrewSupport}
\usepackage{../../../../tex/packages/mathShortcuts}
\usepackage{../../../../tex/packages/theoremsSupport}

\usepackage{listings}
\lstset{
	language=Python,
	numbers=left,
	breaklines,
	basicstyle=\ttfamily,
	commentstyle=\color{gray},
	prebreak=\textrightarrow,
}

\newcommand\low{\mathrm{low}}

\author{שחר פרץ}
\title{אלגוריתמים $\sim$ \textit{הרצאה 3}}
\begin{document}
	\maketitle
	
	בשיעורי הבית מותר להשתמש בחומרים של התרגול בפקולטה. 
	
	\defi{בגרף לא מכוון קודקוד מנתק הוא קודקוש הסרתו מהגרף מנתקת את רכיב הקשירות שלו. }
	\rmark{כשמורידים קודקוד מסיר את כל הקשתות שלו}
	
	
	\textit{אלגוריתם למציאת קודקודים מנתקים: }
	נגדיר לכל $v \in V$ ערך $\low[v]$, שנסמן כמה גבוה בעץ ה־DFS של $v$ ניתן לטפס מצאצא של $v$, או מ־$v$ עצמו (ע''י קשת אחורית). משתמשים בערכי הגילוי של הקודקודים. 
	
	\fig{\begin{forest}
			[1 [2, label=$w$, name=two[3, label=$v$[4, label=$u$, name=four], [5, label=$x$]]], [7[8][9]]]
			\path (two) -- (four);
	\end{forest}}{}
	דמיינו קשת בין 2־ ל־4 שאני אוסיף אחכ (מחוץ ל־DFS). אז $\low[w] = 2$, $\low[u] =2, \low[v] = 2, \low[x] = 5$. 
	
	משתמשים בערכי הגילוי של הקודקודים לצורך השוואה. מעדכנים את $\low$ תוך כדי טיפול בקודקוד, ובחזרה מהשכנים שלו. אם $(u, v)$ קשת אחורית, אז $\low[u] \le \dd[v]$. אם $w$ צאצא של $u$, אז $\low[u] \le \low[w]$, תמיד $\low[u] \le \low[d]$. בזמן סיום טיפול בקודקוד, נתקלנו בקשתות האחריות ממנו, וסיימנו לטפל בילדים שלו. ולכן נוכל לחשב את ה־$\low$ שלו. 
	
	\theo{קודקוש $v$ שאיננו שורש הוא מנותק אמ''מ קיים $w$ בן של $v$ כך ש־$\low[w] \ge \dd[v]$ (כלומר, $v$ מנתק עבור בן שלו $w$ לא נוכל לטפעל מעל $v$). עבור שורש, הוא מנתק אמ''מ יש לו לפחות שני בנים. }\begin{proof}
		עבור שורש – הטענה טרוויאלית. 
		
		עבור קשת שאיננה שורש, אם קיים $w$ בן של $v$ כך ש־$\low[w] \ge \dd[v]$, לא ניתן לטפס מ־$w$ לאבא של $v$ ולכן $v$ מנתק. 
		
		מצד שני, עבור אם $v$ מנתק, כלומר קיים בן שלו $w$ שלא ניתן לטפס ממנו מעל $v$, ולכן $\low[w] \ge \dd[u]$. 
	\end{proof}
	
	נחזור שוב על מושג של $\low$ – ''כמה ניתן לטפס בתת־העץ ב־DFS``. 
	
	\subsection*{רכיבי קשירות חזרה (רק''חים, Strong connected compoenents, SCC)}
	שאלה – איך נוכל לחלק גרף ל־SVV־ים (תת גרף מקסימלי שהוא קשור חזק, כלומר יש מסלול בין כל שני קודקודי ראה הגדרהם)? 
	
	\defi{$G$ קשיר חזק אמ''מ לכל זוג צמתים $u, v$, יש מסלול $u \to v$ ומסלול $v \to u$}. 
	\defi{$G$ מתפרק לרקח''ים, כל אחד מהם מורכב מקוצה מקסימלית של צמתים, שיחד עם הרשות המחברות בינהן, מהוות גרף קשירו חזק. }
	
	לא ייתכנו שני רקח''ים עם צומת משותף. 
	
	אבחנה – רכיב קשירות חזק נמצא באותו העץ ב־DFS (עקרון המסלול הלבן). זה הכרחי, אך לא מספיק. תנאי מספיק יהיה שהוא יהיה לא רק בעץ ה־DFS, אלא שהוא יהיה גם בעץ ה־DFS של ה־transpose (הגרף עם הרשתות ההפוכות). ההגיון פשוט – בהינתן איזשהו s, לא רק שצריך שיהיו קשתות ממנו לשאר האיברים ברכיב הקשירות החזקה (מה שמוודאים כאשר מריצים DFS ומקבלים שהם באותו העץ), אלא שגם יש קשת בינהם אליו (את זה מקבלים כאשר מריצים על ה־transpose). 
	
	זה מוביל אותנו לאלגוריתם Kosaraju-Sharir (1981). נתבונן ב־$G = (V, E)$, ונגדיר $E^{T} = \{(v, u) \mid (u, v) \in E\}$. האלגוריתם: 
	\begin{enumerate}
		\item נריץ DFS על $G$
		\item נסדר את צמתי $G$ לפי סדר $f$ יורד, נסמנם ב־$L^{T}$
		\item נבנה את $G^{T}$. 
		\item נריץ DFS על $G^{T}$, כאשר בלולאה הראשית נעבור לפי הסדר ב־$L^{T}$. 
		\item כל עץ ביער ה־DFS השני, הוא רק''ח. 
	\end{enumerate}
	
	\theo{אם $C$ רק''ח של $G$ אזי כל קודקודי $C$ נמצאים בעץ אחד של ה־DFS השני (וגם הראשון). }
	הוכחה ישירות מעקרון המסלול הלבן, והעובדה שברק''ח יש מסלול בשנ הכיוונים בין זוג קודקודים. 
	\defi{\textit{גרף העל של $G$} יסומן $H(G)$ או $G^{\mathrm{SCC}}$, הוא גרף שצמתיו הם הרק''חים של $G$, וקשת מכוונת מחברת בין שני רק''חים $C_1$ ל־$C_2$ אם קיימת קשת ב־$G$ בין קודקוד $u \in C_1$ ו־$v \in C_2$. קשתות של $H(G)$ נקראות \textit{קשתות על}. }
	
	\theo{$H(G)$ הוא גרף מכוון חסר מעגלים (אנ': אציקלי, DAG). }
	מעגל ב־$H(G)$ היה מחבר כמה ק''חים יחד. 
	\theo{עבור $C_1, C_2$ הק''חים של $G$ וקשת $C_1 \to C_2$, מבין צמתי $C_1 \cup C_2$ הצומת על זמן הסיו ם$f$ המקסימלי מהשלב הראשון נמצא ב־$C_1$. }\begin{proof}
		נסמן ב־$u$ את הקודקוד הראשון מבין $C_1 \cup C_2$ שהתגלה. נפרק למקרים:
		\begin{itemize}
			\item אם $u \in C_1$, יש מסלול לבן ממנו ל־$C_1 \cup C_2$, ולכן הוא יהיה עם זמן הסיום המקסימלי כנדרש. 
			\item אם $u \in C_2$, נחסה את כל $C_2$ לפני שנראה קודקוד כלשהו מ־$C_1$ (אחרת הם חלק מאותו רכיב קשירות, ראה משפט לעיל) ולכן זמן הסיום גבוה יותר כנדרש. 
		\end{itemize}\envendproof
	\end{proof}
	
	נשאר להראות שכל עץ ב־DFS השני (על $G^{T}$) מורכב מרק''ח בודד (הסברנו למה רק''ח תמיד נכנס בשלמותו לעץ). זה ישלים את הוכחת נכונות האלגוריתם. 
	\begin{proof}
		נוכיח באינדוקציה על סדר בניית העצים ב־DFS השני. העץ הראשון, מתחיל בשורש $r$ (הקודקוד עם זמן סיום מקסימלי). נסמן ב־$C_1$ את הרק''ח של $r$. אז כל איברי $C_1$ יופיע בעץ הראשון. נניח בשלילה שרק''ח נוסף $C_2$ יופיע באותו העץ. בה''כ נניח ש־$C_2$ הרק''ח הראשון שנדחף לעץ למרות שממש לא הוזמן (אפשר לקחת מינימום על ריצות ה־DFS), כלומר קיימת קשת מ־$C_1$ ל־$C_2$ ב־$G^{T}$. מכאן שב־$G$ יש קשת על בין $C_2$ לבין $C_1$. לפי הטענה ממקודם מבין $C_1, C_2$ הקודקוד עם זמן הסיום המקסימלי הוא ב־$C_2$, בסתירה לזה שהתחלנו מ־$r \in C_1$. 
		
		הוכחת צעד האינדוקציה זהה להוכחת צעד הבסיס. 
	\end{proof}
	
	
	נמשיך לאלגוריתם נוסף לחישוב או למציאת רק''חים במעבר DFS יחיד (Tarjan, 1972)
	
	הרעיון: נחזיק במחסנית צמתים לפי סדר גילוי הצמתים. נרצה שברגע שמסיימים לטפל ברק''ח להוציא את הקודקודים שלו מהמחסנית. נרצה לטעון, שכל הקודקודים של הרק''ח הנוכחי, רציפים במחסנית.  
נקרא לקודקוד הראשון שמתגלה ברק''ח מסויים ''המנהיג שלו	``, ויהיה לו (פרט להמון כוח, כסף, כבוד) ויהיה לו זמן גילוי מינימלי וזמן סיום מקסימלי מבין קודקודי הרכיב. 

	''כשאני מסיים לטפל בקודקוד, אני צריך לבדוק האם הוא מנהיג של רקח`` – משפט שהיה נשמע ממש מוזר בכל קונטקסט אחר. 
	
	עתה נכתוב מסודר את האלגוריתם.
	
	נשנה את DFS: 
	\begin{enumerate}
		\item נחזיק מכסנית $S$ שנכניס אליה בתחילת DFS-visit
		\item נתחזק $\low[u]$ לכל קודקוד $u \in V$. נאתחל $\low[u] = \dd[u]$ בתחילת הטיפול ב־$u$. 
		\item במעבר על הקשתות שיוצאות מ־$u$ ב־DFS-visit, עבור קשת $(u, v)$ ולאחר טיפול ב־$v$ (אם צריך), במידה $v \in S$ נעדכן $\low[u] = \min \{\low[u], \low[v]\}$. בסיום טיפול ב־$u$, במידה ומתקיים $\low[u] = \dd[u]$, $u$ מנהיג של הרק''ח. נוציא מהמחסנית איברים עד שנגיע ל־$u$ כולל ונסיים רק''ח. 
	\end{enumerate}
	
	אבחנה: $C_1, C_2 \dots C_j$ רקח''ים שהמנהיגים שלהם המנהיגים של הרכיבים האלה גם על המסלול האפור. במחסנית $S$ הקודקודים שגילינו של כל רכיב $C_i$, ברצף, לפי הסדר (המנהיג הוא הראשון מבינהם). 
	
	\theo{$u$ מנהיג של רכיב אמ''מ כשה־DFS נסוג מ־$u$ אין קשת אחורית/חוצה מתת־העץ של $u$ אל מחוץ לתת־העץ של הקודקוד שב־$S$. }
	נרצה לטעון ששומרים על הטענות הבאות: 
	\begin{itemize}
		\item הרכיבים שכבר גילינו נכונים
		\item אם נסוגנו מ־$v$ והוא לא שורש, אז הוא באותו הרכיב כמו ההורה שלו. 
	\end{itemize}
	
	\begin{proof}
		באינדוקציה על צעדי האלגוריתם. הבסיס נכון באופן טרוויאלי. בצעד האינדוקציה, אם $v$ לא מנהיג לפי האלגוריתם, ניתן להגיע ממנו לקודקוד שב־$S$ עם ערך $d$ גילוי נמוך יותר, ולכן לא מנהיג. אם מתקיים ש־$\low[v] = \dd[v]$, כל הקודקודים מ־$v$ עד סוף המחסנית בתת־העץ שלו באותו רכיב כמוהו ומהווים רק''ח. 
	\end{proof}
	
	
	באופן כללי, הדבר האינטואיטיבי הראשון לעשות כשיש לנו גרף מכוון זה לפרק אותו לרכיבי קשירות חזקה ואז לעשות דברים על גרף העל. 
	
	שאלה: בהינתן $G = (V, E)$ גרף לא מכוון, $S \subset V$, רוצים לבדוק האם קיים מסלול שעובר דרך כל קודקודי $S$. המסלול לאו דווקא פשוט, ויכול לעבור גם בקודקודים שאינם ב־$S$. 
	
	באופן כללי, כאשר עוברים לעל־גרף אנחנו ב־DAG ואז אפשר לבצע מיון טופולוגי. זה יעזור לנו כאן: 
	\begin{itemize}
		\item נמצא את $H(G)$
		\item נמצא מיון טופולוגי לגרף העל
		\item עתה, יש סדר לרכיבי הקשירות החזקה ש־$S$ מופיעים בהם. מכאן, שאם קיים מסלול כזה, הוא צריך להתחיל ברק''ח הראשון במיון, ולהמשיך לפי הסדר שנוגע ב־$S$ (לא נוכל ללכת אחורה במיון הטופולוגי, זה ליטרלי בלתי אפשרי). 
	\end{itemize}
	הכל לינארי. 
	
	עתה, נבנה את הפתרון מסודר. 
	\begin{proof}[פתרון]
		נבנה את גרף העל $H(G)$. נבצע מיון טופולוגי שלו. יהיו $C_1, C_2 \dots C_k$ הרק''חים של $G$ שנחתכים עם $S$, לפי סדר המיון. בתוך רק''ח יש מסלול מכוון בין כל זוג קודקודים, ולכן מספק לבדוק אם יש מסלול ב־$H(G)$ ב־$C_{i + 1} \to C_j$ לכל $i \ge 1 < k$. מחזיר שיש מסלול כנדרש אם קיימים ה־$k - 1$ מסלולים האלו ב־$H(G)$. 
	\end{proof}
	
	אנחנו אומנם עומדים להריץ $k - 1$ BFS־ים, אבל אנחנו עומדים להתעלם מקשתות שממשיכות לפני הצומת שנמצאת ב־$S$ הקרובה (כלומר נשתמש רק בקשתות בין $C_i$ לבין $C_{i + 1}$). מכאן שכל קודקוד משתתף בלכל היותר שני BFS־ים, וכל קשת ב־BFS אחד (והחסם לינארי). 
	
	
	נעבור לשאלה אחרת. יהא $G = (V, E)$ גרף מכוון. נרצה למצוא קבוצה מינימלית בגודלה $S$ כך שלכל $v \in V$ קיים $u \in S$ עבורו יש מסלול מ־$u$ ל־$v$. 
	
	נבחר קודקוד אחד (נציג) עם דרגת כניסה $0$ מכל רק''ח. זמן ריצה לינארי בבירור. 
	
	\begin{proof}[נכונות]
		נראה שלכל $v \in V$ קיים $u \in S$ עם מסלול מ־$u$ ל־$v$. יהיו $C_1, C_2, \dots C_k$ רק''חים לפי סדר מיון טופולוגי. נוכיח באינדוקציה על $i$ בעבור $v \in C_i$. בסיס $v \in C_1$, אז ל־$C_1$ דרגת כניסה $0$, ולכן קיים $u \in C_1 \cap S$ שממנו מסלול ל־$v$. 
		
		צעד: נניח נכונות עד $C_{i -1}$. או דרגת הכניסה של $C_i$ היא $0$, כמו הבסיס, אחרת קיימת קשת $C_j \to C_i$ עבור $j < i$ ומה.א. ניתן להגיע לכל קודקודי $C_j$ ולכן לקודקודי $C_i$ [מילה שאני לא מבין] $v$. 
		
		נוכיח ש־$\sof{S}$ מינימלית. אם הייתה $\sof{S'} \le \sof{S}$ המקיימת את הנדרש, קיים $c$ עם דרגת כניסה $0$ בלי מציג ב־$S$, ומכאן לא ניתן להגיע לקודקודיו ב־$C$ רכיב הקשירות שלו. 
	\end{proof}
	
	
	שאלה נוספת להנאתנו. 
	נתון $G = (V, E)$ גרף מכוון אציקלי. צריך למצוא מסלול ארוך ביותר ב־$G$. נציג שני פתרונות: 
	;\begin{itemize}
		\item פתרון אייל: להריץ DFS כך שהלולאה החיצונית \textit{והשכנים} (הלולאה הפנימית) ממוינים לפי הסידור הטופולוגי. אבל, צריך לנתח את עבודת ההכנה הזאת. 
		\item נמצא מיון טופולוגי $v_1, \dots v_n$. נחשב לכל $i$ מהסוף את $\ml_i$, אורך המסלול הארוך ביותר שמתחיל ב־$v_i$ (שימו לב שזה $\oc(n)$, כי אפשר לשרשר ברגע שמגיעים למסלול מוכר). כלומר, נגדיר: 
		\[ \ml_i = \begin{cases}
			0 & \text{אין קשתות}
			1 + \max \{\ml_j \mid (v_i, v_j) \in E\}
		\end{cases} \]
	\end{itemize}
	
	\ndoc
\end{document}